% !TEX program = xelatex
% 2026年11月绍兴市高考科目诊断性考试 语文参考答案与评分标准
\documentclass[11pt,a4paper]{ctexart}
\usepackage{ywexam}

\examfooter{语文参考答案}

\begin{document}

\answertitle{2026年11月绍兴市高考科目诊断性考试}{语文参考答案与评分标准}

\sechead{一、阅读（73分）}

\subsechead{（一）阅读Ⅰ（19分）}
\ans{1}{D。“需要充分地利用本国的化石燃料”有误，与原文“能源独立绝不能建立在对化石燃料的依赖上”的观点相悖。}
\ans{2}{D。此句意在鼓励能源进口国多多使用“风力”“太阳能”等清洁能源，化石燃料容易受外部局势（禁运封锁）影响。}
\ans{3}{B。材料二的核心观点是：治理方案无法落地的根源在于各国政治和经济利益的博弈，处境相异而立场不同。A项强调公平性；C项强调气候问题不是单一危机，而会影响全局；D项展现了应对的积极姿态。故选B。}
\ans{4}{A：气候危机加剧，临界点正在逼近，后果严重，化石燃料是深层风险。（2分）}
\subans{B}{：可再生能源革命已经全面展开，优势明显，这是一场清洁电力的转型（电气化、成本下降、蕴藏着广阔机遇）。（2分）}
\scorenote{（A处强调“气候危机”，B处强调“能源转型出路”，概括原文作答，意思相近即可。共4分。）}
\ans{5}{\xq{1}加快向清洁能源转型，严格限制过冲幅度和持续时间，将全球气温升幅限制在1.5℃内。}
\subans{\xq{2}}{在应对能源危机的同时，减少对化石燃料的依赖，确保能源安全。}
\subans{\xq{3}}{团结协作，消除贫困，改善民生。实现公平转型，发达国家要切实提供更多资金和技术支持。}
\scorenote{（意思相近即可，每点2分，共6分。）}

\subsechead{（二）阅读Ⅱ（18分）}
\ans{6}{C。“虚写”错误。此处与青年人进城无关。}
\ans{7}{B。强调的不是“条件简陋”，而是电影营造的“陌生美好”，与我们不同的世界。}
\ans{8}{\xq{1}表层：因离乡多年，物是人非，还乡时已与乡亲生疏，昔日的亲密不再。}
\subans{\xq{2}}{深层：乡村原有的精神纽带消解，作者即便身在故乡，也再难融入那个曾经温暖热闹的共同体，故而产生精神疏离之感。}
\subans{\xq{3}}{以“故乡”与“异乡人”的矛盾组合，含蓄点明全文的主旨，寄寓了作者对人际温暖和身份归属的深切怀念，有无可挽回的怅惘。}
\scorenote{（意思相近即可，每点2分，共6分。）}
\ans{9}{露天电影，曾是中国乡村常见的社会文化生活与民间娱乐方式。定期在村子之间播放，遇到婚丧嫁娶，村民也常包场以谢邻里。一般在村中空地竖上高杆，中间悬挂白帆布银幕，用放映机投放，男女老少齐聚观看。露天电影不仅丰富了乡村生活，更在集体观影中凝聚了邻里乡情，成为一代人共同的精神家园与文化记忆。}
\scorenote{（来源、特点、文化价值，意思相近即可，各2分，共6分；语言表达不合要求，则酌情扣分。）}

\subsechead{（三）阅读Ⅲ（21分）}
\ans{10}{ADH}
\ans{11}{A。“遽”意为惶恐。}
\ans{12}{A。“两败俱伤”错误，材料一中秦始皇担心的是伤害百姓。句中也无明显因果关系。}
\ans{13}{（1）元年起与百姓一道除旧布新，替他们解除了全部旧有罪责。}
\scorenote{（“黔首”解释为“百姓”；“更始”解释为“重新开始”“除旧布新”；“为（之）”解释为“替他们”；“故罪”解释为“原有罪责”。意译为“大赦天下”等酌情给分，每点1分，共4分。）}
\subans{（2）}{终究不能使它们成为不可更改的经典，只能充作后辈作者删改文稿时的凭借（依据）罢了。}
\scorenote{（“终”解释为“终究”1分，“刊”解释为“更改”1分，“徒”解释为“只能”“白白的”1分，“资”解释为“凭借”“依据”1分，每点1分，共4分。）}
\ans{14}{\xq{1}《赵正书》记载胡亥即位是得到秦始皇认可的，《史记》记载秦二世是靠伪造诏书即位的。\xq{2}材料一透露赵正是主动安排后事，材料三显示生前没有人谈及。\xq{3}材料一透露胡亥即位后赐死扶苏等，材料三显示假托遗诏赐死扶苏后即位。\xq{4}材料一透露杀人是胡亥的决策，赵高并未参与；材料三显示李斯赵高他们参与了决策。}
\scorenote{（任选3条作答，前2条各1分，答对3条给4分。另有言之成理的也酌情给分。）}

\subsechead{（四）阅读Ⅳ（9分）}
\ans{15}{C。A项考查起句作用；B项考查意象与白描手法；D项考查以自然永恒反衬人事沧桑。C项暗考虚实关系。结合词句中的“问”字以及“吹箫人在何许”，可知昔日扬州“柳津花渡”“露桥夜月”下的繁华，是词人追忆、想象（虚写）的场景，是借用杜牧诗意引发的发问，并非词人此刻“极目所见”的真实画面。C项将“虚景”误认作“眼前实景”，故选C。}
\ans{16}{\xq{1}上阕写众人谈论昔日胜景，却只能“谩怀古”，“谩”是徒然、空自之意，写出无力改变现实的深重无奈与悲凉。（2分）}
\subans{\xq{2}}{下阕以“问柳津花渡”领起，追寻昔日繁华，得到的却是“粉黛成黄土”的落空。（2分）}
\subans{\xq{3}}{结尾进一步追问废墟之上的繁华是否还能重现。这一问，将情感从对往昔的追忆转向对未来的深切忧虑与迷茫，与开头的“谩”遥相呼应，勾勒出由无奈到追寻、再由追寻跌入深沉忧思的完整情感脉络。（语意对即可，每点2分，共6分。）}

\subsechead{（五）名篇名句默写（本题共1小题，6分）}
\ans{17}{（1）又闻子规啼夜月，愁空山}
\subans{（2）}{铁衣远戍辛勤久，玉箸应啼别离后}
\subans{（3）}{而智勇多困于所溺，岂独伶人也哉}
\scorenote{（每空1分，共6分。）}

\sechead{二、语言文字运用（17分）}
\ans{18}{B。“泼辣”不准确。}
\ans{19}{A：无伤大雅／无关紧要／无足轻重（能表示小疏漏“不严重、不碍事”，与“大祸将至”相对，搭配得当即可。2分）}
\subans{B}{：万念俱灰／心灰意冷（能表示希望破灭后的极度失望、绝望，2分）}
\ans{20}{\xq{3}改为“便很难精准（地）把握日常表达里细腻的分寸感”（2分）}
\subans{\xq{5}}{改为“将导致中文语感的损害甚至缺失”（2分）}
\ans{21}{\xq{1}汉语重“意合”，语义深藏于场景之中，AI可能难以理解语境。（无法识别其以戏仿、反差营造荒诞幽默的使用场景）}
\subans{\xq{2}}{它仅凭词汇概率推演，可能受训练数据偏见影响。（网络热梗、恶搞文本大量存在，AI可能将之当作常见搭配甚至事实）}
\subans{\xq{3}}{没有中国文化的正确认知，无法理解倒拔垂杨柳的形象和林黛玉的病弱形成反差式趣味。（每点2分，共6分，意思相近即可。）}

\sechead{三、写作（60分）}
\ans{22}{（60分）作文阅卷参照高考评分要求。}

\par\vspace{0.5em}
\sechead{文言文参考译文}

\translabel{材料一：}
\begin{translation}
从前，秦王赵正巡游天下，返回途中到达柏人时生了病，病情危重。他长叹着流泪，深深地叹息，随后召来丞相李斯，告诉他说：“我自己审度天命，命中注定五十岁而死。我活到十四岁就被立为王，至今在位已经三十七年了。我本当在今年死去，却不知道具体是哪月哪日，所以才巡游天下，想借此改换气数、更改天命——难道这样做竟不行吗？如今病重，快要死了。你务必要日夜兼程，赶到甘泉的驿站，不要等待后面的车队。要小心谨慎地处理此事，千万不要让群臣知道我生病了。”

秦王病势随即大大加重，已经不能继续前行，所以再次召来丞相李斯，说：“我作为霸王所享的寿数已经足够了，不知我该拿孤弱年幼的儿子们怎么办啊。我听说过这样的话：牛马相斗，蚊虻就死在它们蹄下；大臣相争，平民百姓就要受苦。我哀怜我孤弱的儿子以及我抚育的臣民，这份挂念到死也不会忘记。你们商议一下立谁为继承人吧。”丞相臣李斯、御史臣冯去疾冒死叩头进言说：“如今返回的路途遥远，而若回京后再定期限召集群臣，恐怕大臣中会有人图谋不轨，请求立公子胡亥为王位继承人。”秦王说：“可以。”

秦王死后，胡亥即位，随即杀死了他的兄长扶苏和中尉蒙恬。他大赦罪人，又把原为隶臣刑徒的赵高赦免出来，让他担任郎中令。接着诛灭那些王室宗族，毁坏了国家（土谷神庙），焚烧那些律令以及前代典藏。他还想动用万乘车队的仪制来镇抚天下，说：“将要与天下人一道重新开始，除旧布新。”
\end{translation}

\translabel{材料二：}
\begin{translation}
天下失去始皇帝，（天下人）都惊惧悲哀至极。我（胡亥）是奉行先帝遗诏继位。如今，宗庙之事以及著于简册、用来彰明至治大功德（始皇帝的丰功伟绩）的（典章）都已完备，应当废除和修订的律令也都处理完毕。（朕将以）元年（为起点）与百姓一道除旧布新，把旧有的罪责全部解除。（相关的）法令都已颁布了，我将亲自镇抚天下。官吏、百姓要各司其职、认真办事，不要用徭役赋税烦扰百姓，也不要因琐碎小事苛细举劾县吏。要尽快将此诏公布。
\end{translation}

\translabel{材料三：}
\begin{translation}
（始皇一行）到达平原津就病了。始皇厌恶别人说死，群臣没有谁敢提死后的事情。皇上的病越发加重，于是就写下加盖玉玺的诏书赐给公子扶苏说：“（你）回咸阳来参与治办丧事，会合（灵柩）安葬（朕）。”诏书已经封好，放在中车府令赵高兼管符玺事务的官署，还没有交给使者。七月丙寅日，始皇在沙丘平台驾崩。丞相李斯因为皇上驾崩在都城之外，担心诸位公子和天下发生变乱，就隐瞒了这件事，不发布丧讯。把棺材装载在辐凉车中，让原先受宠的宦官陪乘，车驾所到之处照旧进献食物。百官像往常一样奏报政事，宦官就从辐凉车中批准他们的奏事。唯独皇子胡亥、赵高和（始皇）宠幸的宦官五六人知道皇上已死。

赵高原先曾教胡亥书法和刑狱法令之事，胡亥私下亲近他。赵高就与公子胡亥、丞相李斯暗中合谋，毁掉始皇封好赐给公子扶苏的那封诏书，而另外伪造了一份丞相李斯在沙丘接受的始皇遗诏，立皇子胡亥为太子。还写下一份矫诏赐给公子扶苏、蒙恬，列举（他们的）罪状，（命他们）自裁。（车队）经直道到达咸阳，才发布丧事。太子胡亥继承帝位，就是二世皇帝。
\end{translation}

\translabel{材料四：}
\begin{translation}
大体说来，偏纪、小录这一类书，记的都是当天当时发生的事，同国史相比，它们算是（记载）最为真实的了。然而它们的言辞大多粗鄙质朴，记事很少完备周全，终究不能使它们成为不可更改的经典，只能充作后辈作者删改文稿时的依据罢了。所谓逸事，都是前代史书遗漏未载、后人记录下来的，（学者）从各种不同的说法中去寻访，（它们）带来的益处确实很多。可是到了虚妄之人编写这类书时，就随意记载传闻，而不加权衡选择，于是真伪不分、是非相乱——这是它弊端最严重的地方。
\end{translation}

\end{document}
